\chapter{Problèmes et graphiques}\label{ch:g5:data}

Le dernier chapitre de l'année met tout au travail~: des problèmes
à plusieurs étapes avec de l'argent, des quantités et des durées, et
les graphiques qui présentent l'information d'un coup d'œil.
L'idée phare est le «~par un~» --- trouver le prix d'un seul
article --- qui grandira en proportionnalité dans
\cref{ch:g6:propdata}.

\section{Problèmes à plusieurs étapes}

\begin{method}[Apprivoiser un long problème]\label{met:g5:data:steps}
\begin{enumerate}
  \item Lire jusqu'au bout~; dire quelle est la question finale~;
  \item planifier à rebours~: pour y répondre, de quoi ai-je besoin
        d'abord~?
  \item résoudre les étapes dans l'ordre, en écrivant une opération
        et une phrase courte par étape~;
  \item vérifier la réponse finale par le bon sens et par une
        estimation.
\end{enumerate}
\end{method}

\begin{example}\label{ex:g5:data:multistep}
«~Une école commande $4$ cartons de $25$ livres à $6$ chacun, et
paie $180$ de livraison en plus. Quel est le coût total~?~»
\begin{enumerate}
  \item Livres~: $4 \times 25 = 100$ livres.
  \item Coût des livres~: $100 \times 6 = 600$.
  \item Total~: $600 + 180 = 780$.
\end{enumerate}
Vérification par estimation~: environ cent livres à $6$, c'est
environ $600$, plus environ $200$~: environ $800$ --- la réponse
$780$ est plausible.
\end{example}

\begin{example}[Par un]\label{ex:g5:data:perone}
«~$5$ mugs identiques coûtent $12.50$~; que coûtent $8$ mugs~?~»
Trouver d'abord le prix d'\emph{un} mug~:
\[
12.50 \div 5 = 2.50,
\qquad\text{puis}\qquad
8 \times 2.50 = 20 .
\]
Par le un, tout devient accessible --- ce schéma en deux étapes
résout toute une famille de problèmes (recettes, vitesses, prix),
et c'est la graine de \cref{ch:g6:propdata}.
\end{example}

\begin{example}[Monnaie et budget]\label{ex:g5:data:budget}
Zoé a $30$. Elle achète un livre à $12.40$ et deux stylos à
$2.30$ chacun. Peut-elle aussi s'offrir un jeu de société à
$13$~?
\begin{enumerate}
  \item Stylos~: $2 \times 2.30 = 4.60$.
  \item Dépensé jusqu'ici~: $12.40 + 4.60 = 17$.
  \item \omterm{def:g3:division:remainder}{Reste}~: $30 - 17 = 13$ --- exactement assez pour le jeu,
        sans rien laisser.
\end{enumerate}
\end{example}

\section{Lire des graphiques}

\begin{method}[Lire n'importe quel graphique]\label{met:g5:data:read}
\begin{enumerate}
  \item Lire le titre~: que compte-t-on~?
  \item Lire les axes ou la légende~: que représente une
        graduation, une barre, un symbole~?
  \item Seulement ensuite répondre aux questions --- pointer le
        graphique avec une règle aide à lire les hauteurs
        exactement.
\end{enumerate}
\end{method}

\begin{example}\label{ex:g5:data:chart}
Précipitations dans une ville, mois par mois~:

\begin{omfigure}
\begin{tikzpicture}
\begin{axis}[
  width=8.8cm, height=5.0cm,
  ybar, bar width=13pt,
  symbolic x coords={janv.,févr.,mars,avr.,mai,juin},
  xtick=data,
  ymin=0, ymax=90,
  ytick={20,40,60,80},
  ylabel={pluie (mm)},
  tick label style={font=\small},
  label style={font=\small}]
  \addplot[fill=omDef!30, draw=omDef] coordinates
    {(janv.,60) (févr.,45) (mars,55) (avr.,70) (mai,80) (juin,30)};
\end{axis}
\end{tikzpicture}

{\small Un diagramme en barres des précipitations. Une graduation
$= 20$~mm.}
\end{omfigure}

Lectures~: le mois le plus humide est mai ($80$~mm)~; juin a reçu
la \omterm{def:g3:division:half}{moitié} de la pluie d'avril ($30$ contre $70$~? non --- la \omterm{def:g3:division:half}{moitié}
de $70$ est $35$, donc \emph{pas tout à fait} la \omterm{def:g3:division:half}{moitié}~: lire
précisément compte~!)~; total pour les six mois~:
$60 + 45 + 55 + 70 + 80 + 30 = 340$~mm.
\end{example}

\begin{example}[Une quantité qui change]\label{ex:g5:data:line}
La hauteur d'une plante, mesurée chaque semaine~:

\begin{omfigure}
\begin{tikzpicture}
\begin{axis}[omaxis,
  width=8.4cm, height=5.0cm,
  xmin=0, xmax=6.5, ymin=0, ymax=24,
  xlabel={semaine}, ylabel={hauteur (cm)},
  xtick={1,...,6}, ytick={5,10,15,20}]
  \addplot[omDef, thick, mark=*, mark size=1.6pt] coordinates
    {(1,3) (2,6) (3,10) (4,15) (5,19) (6,21)};
\end{axis}
\end{tikzpicture}

{\small Joindre les points de mesure montre la \emph{croissance}~:
plus le segment est raide, plus la plante a grandi cette
semaine-là.}
\end{omfigure}

La plante a grandi le plus vite entre les semaines $3$ et $4$
($+5$~cm) et a ralenti à la fin ($+2$~cm en semaine $6$).
\end{example}

\section{Exercices}

\begin{exercise}[$\star$]\label{exo:g5:data:1}
Résoudre par étapes~: «~Un club loue un bus pour $240$ et achète
$32$ billets de musée à $7$ chacun. Quel est le coût total de la
sortie~?~»
\end{exercise}

\begin{exercise}[$\star$]\label{exo:g5:data:2}
Résoudre avec le «~par un~»~: «~$3$~kg de pommes coûtent $7.50$.
Que coûtent $5$~kg~?~»
\end{exercise}

\begin{exercise}[$\star$]\label{exo:g5:data:3}
Résoudre~: «~Six amis partagent équitablement le coût de $87$ d'un
pique-nique. Combien chacun paie-t-il~?~»
\end{exercise}

\begin{exercise}[$\star$]\label{exo:g5:data:4}
Sam a $25$. Il achète un T-shirt à $13.90$ et une casquette à
$8.50$. Combien d'argent lui reste-t-il~?
\end{exercise}

\begin{exercise}[$\star$]\label{exo:g5:data:5}
Avec le graphique des précipitations de \cref{ex:g5:data:chart}~:
quels mois ont reçu plus de $50$~mm~? Combien de pluie en plus est
tombée en mai qu'en juin~?
\end{exercise}

\begin{exercise}[$\star$]\label{exo:g5:data:6}
Avec le graphique de la plante de \cref{ex:g5:data:line}~: quelle
hauteur avait la plante en semaine $2$~? Entre quelles semaines
a-t-elle grandi d'exactement $4$~cm~?
\end{exercise}

\begin{exercise}[$\star$]\label{exo:g5:data:7}
Dessiner un diagramme en barres pour le nombre de buts marqués par
une équipe~: septembre $4$, octobre $7$, novembre $3$, décembre
$6$. Choisir la graduation, et donner le total.
\end{exercise}

\begin{exercise}[$\star$]\label{exo:g5:data:8}
Un pack de $8$ yaourts coûte $3.20$~; vendu seul, un yaourt coûte
$0.50$. Calculer le prix «~par un~» dans le pack, et combien on
économise par yaourt en achetant le pack.
\end{exercise}

\begin{exercise}[$\star\star$]\label{exo:g5:data:9}
«~Un cinéma a vendu $148$ billets à $9.50$ le samedi et $95$
billets au même prix le dimanche. Combien d'argent le samedi a-t-il
rapporté de plus que le dimanche~?~» Résoudre de deux façons~: en
calculant les recettes des deux jours, ou en calculant avec la
\emph{\omterm{ex:g1:subtraction:difference}{différence}} de billets --- et vérifier que les réponses
coïncident.
\end{exercise}

\begin{exercise}[$\star\star$]\label{exo:g5:data:10}
Une recette pour $4$ personnes demande $300$~g de riz. Aline
cuisine pour $10$ personnes. Combien de riz lui faut-il~? (Par une
personne d'abord --- décimaux autorisés.)
\end{exercise}

\begin{exercise}[$\star\star$]\label{exo:g5:data:11}
Inventer un problème en trois étapes utilisant de l'argent dont la
réponse finale est $5.50$, écrire sa solution complète dans le
style de \cref{met:g5:data:steps}, et le tester sur un camarade.
\end{exercise}
